\begin{sidewaystable}[p]
\fontsize{7.2pt}{8.4pt}\selectfont
\centering
\caption{\textbf{Filter and label validation map.} The table summarizes how each non-native filter or analytic label was assigned, what validation or audit evidence supports it and where it is used in the analyses. Filtering counts are reported in Supplementary Table~\ref{tab:corpus_filtering_pipeline}; full agreement metrics are reported in Supplementary Tables~\ref{tab:label_validation} and \ref{tab:production_label_validation}.}
\label{tab:filter_classifier_validation}
\setlength{\tabcolsep}{4pt}
\renewcommand{\arraystretch}{1.12}
\begin{tabularx}{\textheight}{>{\raggedright\arraybackslash}p{0.15\textheight}>{\raggedright\arraybackslash}p{0.24\textheight}>{\raggedright\arraybackslash}p{0.25\textheight}>{\raggedright\arraybackslash}X}
\toprule
Variable or filter & How it was assigned & Validation or audit evidence & Role in analysis \\
\midrule
Task ecology: coding versus writing & LLM-assisted semantic task filters selected conversations matching coding-oriented or writing-oriented use after minimum-turn screening. The filters were English-oriented and used a score threshold of 0.70 after keyword prefiltering where applicable. & The filtering ledger records retained counts and selection steps for each corpus (Supplementary Table~\ref{tab:corpus_filtering_pipeline}). In the writing-210 validation artifacts, the conversation-topic field showed high human--LLM agreement (F1=0.973, MCC=0.872, Gwet AC1=0.936; support=210), supporting the topic-label boundary used in the task filter. & Defines the coding- and writing-oriented task settings and the task-ecology covariate used in pooled context models. \\
\addlinespace
User framing: intentional versus unintentional & Conversation-level LLM-assisted \texttt{learning\_intent} annotation marked conversations explicitly oriented toward understanding, exploration or learning versus those without explicit learning-oriented framing. & A human-verified audit reviewed 450 first user turns across all six corpus-by-task settings (75 per setting; 50 positive and 25 negative cases according to the corpus-scale labels). Overall intentional-framing agreement was F1=0.857, MCC=0.677 and Gwet AC1=0.671; complete metrics are reported in Supplementary Table~\ref{tab:production_label_validation}. & Used as the observed user-framing variable in Fig.~\ref{fig:engagement_ecology} and the contextual models; prevalence estimates come from the full corpus-scale labels. \\
\addlinespace
Language filtering & WildChat and LMSYS Chat used English-oriented source or semantic filters. ShareChat received an additional strict-English screen after task filtering, retaining only records with final conversation-level language flag \texttt{English}. & ShareChat strict-English retained 2,481 of 3,494 coding conversations and 1,539 of 1,663 writing conversations after task filtering (Supplementary Table~\ref{tab:corpus_filtering_pipeline}). In the writing-210 validation artifacts, language showed near-perfect agreement (F1=0.998, MCC=0.705, Gwet AC1=0.995; support=210), with MCC lower because nearly all reviewed records were English. & Defines the English-oriented analytic sample; the additional strict-English screen applies only to ShareChat. \\
\addlinespace
Turn-level labels: user engagement and scaffolded support & Corpus-scale annotation used the final codebook, parser checks and post-processing rules. Constructive engagement and scaffolded support are the primary user-side measure and assistant-side measure. & Domain-stratified human--human codebook validation shows strong agreement for constructive engagement (coding F1=0.85; writing F1=0.89) and scaffolded support (coding F1=0.87; writing F1=0.97), with full per-label codebook metrics in Supplementary Table~\ref{tab:label_validation}. A WildChat post-labelling check reviewed 100 labelled turns per task under the observed corpus-scale label distribution (Supplementary Table~\ref{tab:wildchat_post_labelling_review}). Human-confirmed corpus-scale label audits across all six corpus-by-task settings showed acceptable or better reliability for user engagement and assistant labels: the 600-case user audit had constructive-vs-non-constructive F1=0.922, MCC=0.840 and Gwet AC1=0.837; the 180-case assistant audit had scaffolded-support F1=0.912, MCC=0.590 and Gwet AC1=0.791, with support-form F1 values ranging from 0.691 to 0.938 (Supplementary Table~\ref{tab:production_label_validation}). & Provides the turn-level measures used in prevalence, association, support-form and adjacent-turn analyses. \\
\bottomrule
\end{tabularx}
\end{sidewaystable}
